\section{Related Work}

\subsection{Reinforcement Learning from Execution Feedback}

RLEF with binary execution rewards was established by \citet{Gehring2024RLEF}, who demonstrated order-of-magnitude inference-time sample reductions through binary pass/fail rewards combined with GRPO policy optimization. This work confirmed that binary execution correctness is sufficient for meaningful policy improvement and set the standard MBPP training $\to$ HumanEval+ evaluation protocol that we follow.

The mathematical foundation for why binary rewards lead to gradient starvation is the GRPO standard deviation identity: $\sigma = \sqrt{k(G-k)}/G$, where $k$ is the number of correct completions out of $G$ \citep{BayYearickLab2025GRPOIdentity}. When $k=0$ or $k=G$, $\sigma=0$ and no gradient is produced.

\subsection{Gradient Starvation in GRPO}

The empirical consequences of gradient starvation were quantified by \citet{Nie2026GradientStarvation}: at $G=4$ with binary rewards, 54.75\% of training groups have all-fail outcomes and 14.50\% have all-pass outcomes, producing 69.25\% zero-gradient groups in expectation. Our work provides, to our knowledge, the first \textit{per-problem} systematic empirical characterization of this effect for MBPP under constrained generation, finding an even more severe all-fail rate (91.7\%).

\subsection{Online Data Selection for GRPO Efficiency}

Multiple recent papers converge on selecting intermediate-difficulty problems as the key to GRPO efficiency --- all as \textit{online} methods.

\citet{Jiang2026VIGOR} allocate rollouts based on variance-utility signals computed during active training, paralleling our offline profiling conceptually but requiring training loop integration.

\citet{Baroian2026PromptReplay} track per-problem running pass rates during training and prioritize problems with pass rate $\approx 0.5$ --- the maximum-variance point under the binary variance identity. This is the closest conceptual predecessor to our work, but cannot operate as a pure offline step.

\citet{Sun2025DataEfficiency} demonstrate difficulty-targeted online selection for math reasoning with scalar rewards, achieving 23--62\% compute reduction. They target a different reward type and domain; we extend the data selection intuition to binary-reward code generation with an offline approach.

\citet{Cui2026LZE} fuse pass-rate momentum and outcome-uncertainty during training. Like the above, this requires an active training loop.

\textbf{Our approach} occupies the offline niche: profile once with a frozen model, train anywhere. The cold-start finding we identify --- that parameter mismatch between profiling and training can nullify the entire selection --- is the key contribution that prior online methods naturally avoid (they compute difficulty during training, so profiling-training alignment is not a concern).

\subsection{Positioning}

We contribute: (1) the first offline variance-guided selection method for binary-reward code generation RLEF, (2) to our knowledge, the first systematic per-problem binary reward variance distribution characterization on MBPP (prior GRPO-on-MBPP papers encounter this distribution but do not report it at per-problem granularity), and (3) the first systematic identification of cold-start as a precondition for data selection in RLEF.
